% Generated draft: requires mathematical and source review.
\documentclass[11pt]{article}
\usepackage[margin=25mm]{geometry}
\usepackage{fontspec}
\setmainfont{Latin Modern Roman}
\usepackage{amsmath,amssymb,mathtools,unicode-math}
\setmathfont{Latin Modern Math}
\usepackage{enumitem,needspace,xurl,hyperref}
\setlength{\parindent}{0pt}
\setlength{\parskip}{5pt}
\setlength{\emergencystretch}{3em}
\newcommand{\N}{\mathbb{N}}
\newcommand{\Z}{\mathbb{Z}}
\newcommand{\Q}{\mathbb{Q}}
\newcommand{\R}{\mathbb{R}}
\newcommand{\C}{\mathbb{C}}
\newcommand{\F}{\mathbb{F}}
\newcommand{\A}{\mathbb{A}}
\newcommand{\PP}{\mathbb{P}}
\newcommand{\E}{\mathbb{E}}
\newcommand{\eps}{\varepsilon}
\newcommand{\dd}{\,\mathrm{d}}
\begin{document}
\section*{UnsolvedMath ID: 30003050}
\textbf{Source problem code:} catalogue entry
\subsection{Definitions and mathematical statement}
A switching system consists of a collection of evolution laws together with a class of admissible switching signals specifying which law acts at each time. The supplied record does not specify the time domain, switching constraints, state space, or solution concept.

For notation only, a discrete-time realization on a state space \(X\), with modes indexed by a set \(I\), has the form
\[
 x_{k+1}=f_{\sigma(k)}(x_k),\qquad k\in\{0,1,2,\ldots\},
\]
where \(f_i:X\to X\) and \(\sigma\) belongs to a specified class of admissible switching signals. A switched linear realization has \(f_i(x)=A_i x\) on a vector space. This notation is illustrative; the source does not restrict the question to discrete time or unrestricted switching.

Given a complex vector space \(\mathcal F\) of functions \(g:X\to\C\) that is invariant under composition with every \(f_i\), the modewise Koopman operators are the linear maps
\[
 U_i:\mathcal F\to\mathcal F,\qquad (U_i g)(x)=g(f_i(x)).
\]
If \(\mathcal F\) is infinite dimensional, these are infinite-dimensional linear operators. Their boundedness, domains in an unbounded-operator formulation, and spectral properties require additional choices not supplied by the record.

\textbf{Source question.} Can the joint spectral quantities and the stabilizability radius of switched linear systems be extended to suitable classes of infinite-dimensional linear operators, including Koopman operators? Can one associate a canonical Koopman operator with a given switching system?

\textbf{Unresolved formulation.} The record does not identify the particular joint spectral quantities, define its stabilizability radius, or specify the properties an extension must preserve beyond agreement with the intended finite-dimensional notions. It also does not specify whether stabilization uses a single signal, signals selected separately for each initial condition, or a feedback rule. For the second question, the observable space, treatment of switching information, and meaning of canonicality or uniqueness are unspecified. Thus the displayed modewise construction supplies notation, but does not identify the single operator requested by the source. The record supports these research questions, not a fully quantified existence or uniqueness conjecture.
\subsection{Short English statement}
The problem asks whether joint spectral and stabilization measures for switched linear systems extend to infinite-dimensional operators, particularly Koopman operators, and whether a switching system determines a canonical Koopman operator. The operator setting, admissible switching, stabilization criterion, and meaning of canonicality remain unspecified, so a precise theorem-level formulation requires further source clarification.
\subsection{Sources}
\begin{itemize}
\item Supplied local UnsolvedMath registry record: permanent ID \(30003050\), title \emph{Koopman Operators for Switched Infinite-Dimensional Systems}, exact source problem code \texttt{catalogue entry}, in the Oberwolfach Reports Open Problems collection. The supplied proposed year is 2016; no individual proposer is identified.
\item Supplied source citation: \emph{Mini-Workshop: Applied Koopmanism} (2016), Oberwolfach Reports, pp.\ 297--340, DOI \texttt{10.4171/owr/2016/7}. Original report: \url{https://doi.org/10.4171/owr/2016/7}. The background describes Oberwolfach Reports as documenting mathematical workshops and their open-problem discussions; no problem-specific page or named contributor is supplied.
\item Supplied literature source, labeled ``2025 switched-systems Koopman paper'': \url{https://doi.org/10.1016/j.sysconle.2025.106340}. No authors or full bibliographic title are supplied.
\item Literature assessment dated exactly \texttt{2026-08-21}: the supplied curation labels the problem open and states that the proposed infinite-dimensional and unique Koopman construction is not canonical in general, while later switched-system applications do not settle the extracted existence and uniqueness questions. This is supplied dated curation, not an independently verified literature assessment. The background expressly requests independent verification of time-sensitive status and cited claims.
\item The registry labels the statement \texttt{corrected\_verified} and reports a manual semantic, typographic, and LaTeX review asserting preservation of the mathematical target. These are supplied editorial claims, not independent verification in this draft.
\end{itemize}
\end{document}
